\chapter{Diseño y métodos del proyecto}
\label{cap:metodos}

\section{Enfoque y diseño}

El estudio adopta un \textbf{enfoque metodológico cuantitativo}: todo el conocimiento que
produce proviene de medidas numéricas derivadas de las máscaras, y las afirmaciones se
sostienen en estadísticos calculados sobre ellas y no en interpretación de material
cualitativo. Dentro de los tipos de investigación, corresponde a una investigación
\textbf{descriptiva con componente aplicado}: descriptiva porque su propósito es
caracterizar el terreno mediante indicadores y no explicar relaciones causales ni predecir;
y con componente aplicado porque los indicadores se construyen para un uso concreto —la
caracterización del terreno y el examen de lo que las máscaras permiten medir— y no como
ejercicio teórico.

Desde el punto de vista del diseño, el estudio es \textbf{observacional, descriptivo y de
corte transversal}. Observacional
porque no se interviene sobre el objeto de estudio: se analizan máscaras ya producidas, sin
modificar las imágenes ni generar máscaras nuevas. Descriptivo, por las razones ya expuestas.
Y de corte transversal porque cada imagen se analiza por separado, sin modelar su
evolución temporal; el orden de adquisición se emplea únicamente como variable descriptiva y no
como serie temporal calibrada. Las imágenes no son, sin embargo, observaciones independientes:
se adquieren en secuencias, varias en un mismo sol, y tomas próximas pueden mostrar el mismo
terreno. La Sección~\ref{sec:evaluacion} describe cómo se tiene en cuenta esa dependencia en
los contrastes.

La \textbf{unidad de análisis} es la pareja formada por una imagen y su máscara de
segmentación. Cada unidad produce una fila del conjunto de resultados con los indicadores
calculados. No hay agregación previa: las estadísticas del
Capítulo~\ref{cap:resultados} se derivan de esas filas.

Una decisión de diseño conviene explicitarla desde el principio, porque condiciona la
interpretación de todo lo que sigue: \textbf{el procedimiento no lee las imágenes}. Toda
la información que consume proviene de la máscara. Los indicadores E1 y E2 no usan las
imágenes; estas intervienen solo en la inspección visual, en la calibración, en la validación
humana y en los modelos de aprendizaje automático con que se contrastan los indicadores.
Esta elección responde a la pregunta de investigación —qué puede medirse sobre las
máscaras— pero impone un límite que el Capítulo~\ref{cap:resultados} cuantifica.

\section{Datos y fuentes de información}
\label{sec:datos}

\subsection{Origen y versión}

Las fuentes de información de este trabajo son \textbf{secundarias}, con una excepción
acotada. No se toman imágenes nuevas —provienen de una misión espacial— ni se producen
máscaras propias: las generaron terceros. La técnica de recolección es, en consecuencia, la
\textbf{revisión y el procesamiento de un conjunto de datos documental preexistente},
seleccionado según los criterios que se detallan más adelante. La excepción es el conteo por
bandas que el autor recogió para la validación humana (Sección~\ref{sec:evaluacion}), único
dato primario del trabajo.

Se emplea el conjunto AI4Mars en su versión 0.6 \citep{swan2022ai4marsdata}, descargada
del repositorio público. El conjunto reúne imágenes de los rovers \textit{Spirit},
\textit{Opportunity}, \textit{Curiosity} y \textit{Perseverance} procedentes del Planetary
Data System \citep{nasapds2024}, junto con máscaras de segmentación por píxel obtenidas por
anotación colaborativa y, para una parte de las imágenes, máscaras de especialista.

\subsection{Subconjunto de estudio}

El alcance se restringe a la \textbf{cámara de navegación del rover \textit{Curiosity}}
con \textbf{máscaras colaborativas de entrenamiento}. Quedan fuera los rovers de la misión
\textit{Mars Exploration Rovers}, el rover \textit{Perseverance} y la cámara de mástil
multiespectral. La restricción persigue homogeneidad instrumental: todas las escenas
provienen del mismo modelo de cámara —las dos unidades de la cámara de navegación de
\textit{Curiosity}—, con la misma resolución y el mismo régimen radiométrico, de modo que las
comparaciones entre imágenes no arrastran diferencias de cámara.

El emparejamiento entre imagen y máscara se realiza por coincidencia del nombre base. El
directorio de imágenes contiene \cNImagenes{} archivos y el de máscaras de entrenamiento
\cNEscenas{}; el análisis se aplica a las \cNEscenas{} parejas completas, ya que sin máscara
no hay dato de entrada.

Para validar se emplean además las \cNGold{} máscaras de especialista del propio conjunto, que
AI4Mars publica en tres niveles de consenso; salvo indicación, la referencia es el nivel 1.
Conviene advertir desde ahora que las máscaras de especialista \textbf{no recaen sobre las
mismas imágenes} que las colaborativas: ambos conjuntos son disjuntos por construcción. Esta
circunstancia condiciona cómo pueden compararse, y la Sección~\ref{sec:evaluacion} describe el
diseño adoptado para hacerlo.

\subsection{Codificación de las máscaras}

Cada máscara es una imagen en escala de grises del mismo tamaño que la original
—\num{1024}$\times$\num{1024} píxeles— donde el valor entero de cada píxel identifica la
clase de terreno. La Tabla~\ref{tab:codificacion} recoge la codificación de la escala de
navegación.

\begin{table}[htbp]
  \centering
  \caption{Codificación de la escala de navegación en las máscaras de AI4Mars v0.6.}
  \label{tab:codificacion}
  \begin{tabular}{clp{7.6cm}}
    \toprule
    \textbf{Valor} & \textbf{Clase} & \textbf{Descripción} \\
    \midrule
    0   & Suelo (\textit{soil})        & Regolito de grano fino. \\
    1   & Lecho rocoso (\textit{bedrock}) & Roca expuesta en continuidad con el sustrato. \\
    2   & Arena (\textit{sand})        & Depósito de arena suelta. \\
    3   & Roca grande (\textit{big rock}) & Bloque discreto apoyado sobre el terreno. \\
    255 & Sin etiqueta (\textit{null}) & Píxel sin clase asignada. \\
    \bottomrule
  \end{tabular}
\end{table}

Esta codificación se verificó de tres formas independientes antes de calcular ningún
indicador: contra el archivo de claves de clase distribuido con el conjunto, contra su
documentación, y píxel a píxel sobre una muestra de máscaras contrastada con las imágenes
correspondientes. La verificación fue necesaria porque la codificación asumida en la fase
de propuesta era incorrecta —suponía cinco clases con el fondo en el valor cero— y
haberla arrastrado habría invalidado silenciosamente todos los resultados.

El valor \num{255} agrupa dos situaciones distintas: píxeles sin acuerdo suficiente entre
anotadores y píxeles que el propio conjunto excluye por diseño, correspondientes al
cuerpo del rover, al cielo y a distancias superiores a treinta metros. En ambos casos el
píxel se excluye del denominador de la cobertura sobre píxeles etiquetados, y la fracción que
representa se reporta como indicador de calidad de cada escena.

\subsection{Criterios de inclusión}

Cada indicador se define sobre una población explícita, que se mantiene en todo el documento.
Para que no puedan mezclarse poblaciones distintas, la elegibilidad no se deduce de la bandera
de calidad sino que se registra en cuatro columnas del conjunto de resultados, cada una con una
condición explícita —las dos últimas combinan las anteriores—: \texttt{has\_bigrock} (la máscara contiene roca grande),
\texttt{is\_mostly\_null} (más del \SI{95}{\percent} de la imagen sin etiqueta),
\texttt{is\_e1\_eligible} e \texttt{is\_e2\_eligible}. Un único módulo del código define, a
partir de ellas, las dos poblaciones, y todas las cifras de E1 y E2 del documento se obtienen
de ese módulo.

\begin{itemize}
  \item La \textbf{cobertura} (población de E1) se calcula para toda máscara con al menos un
    píxel etiquetado.
  \item El \textbf{conteo} (población de E2) se aplica a las escenas que contienen al menos un
    píxel de roca grande y tienen como máximo el \SI{95}{\percent} de la imagen sin etiquetar, que
    coinciden con las de bandera \texttt{ok}. Se excluyen las \cNNullBig{} escenas que contienen
    roca grande pero están casi enteramente sin etiquetar, porque en ellas la región anotada es
    una fracción ínfima de la escena y el conteo no es interpretable. El número de rocas se
    calcula para todas las escenas, pero solo las elegibles entran en los resúmenes.
\end{itemize}

Ninguna imagen se descarta del conjunto de resultados: las que no admiten un indicador reciben
una bandera de calidad que lo documenta, de modo que el conjunto final describe también sus
propias ausencias.

\section{Métodos y herramientas}

\subsection{Estructura del procedimiento}

La Figura~\ref{fig:esquema-metodo} resume el flujo completo. La máscara de entrada se
bifurca en dos ramas: una calcula la cobertura por conteo de píxeles, y la otra estima el
número de rocas mediante una secuencia de operaciones geométricas.

\begin{figure}[htbp]
  \centering
  % El documento va a doble espacio; dentro del diagrama se fuerza espacio simple
  % para que el texto de los nodos no quede separado.
  \begin{spacing}{1}
  \begin{tikzpicture}[
      font=\footnotesize,
      node distance=5mm and 6mm,
      caja/.style={draw, rounded corners=2pt, align=center, inner sep=3pt,
                   minimum height=9mm, text width=24mm},
      dato/.style={caja, fill=black!8},
      op/.style={caja},
      res/.style={caja, very thick, fill=black!3},
      fl/.style={-{Latex[length=1.8mm]}, thick}
    ]

    \node[dato] (mask) {Máscara\\AI4Mars};

    % --- Rama superior: cobertura (E1) ---
    \node[op, above right=3mm and 9mm of mask] (val) {Roca ($1, 3$) sobre\\válidos ($\neq 255$)};
    \node[res, right=of val] (e1) {\textbf{E1}\\Cobertura};

    % --- Rama inferior: conteo (E2), en dos hileras ---
    \node[op, below right=3mm and 9mm of mask] (bin) {Binaria de\\roca grande};
    \node[op, right=of bin] (morf) {Apertura y\\cierre $3\times3$};
    \node[op, right=of morf] (dist) {Transformada de\\distancia, $\sigma = 3$};

    \node[op, below=8mm of bin]  (ws)   {División\\de aguas};
    \node[op, below=8mm of morf] (filt) {Filtros de\\área y forma};
    \node[res, below=8mm of dist] (e2)  {\textbf{E2}\\Conteo};
    \node[op, left=of ws] (max) {Máximos de\\prominencia};

    \draw[fl] (mask) |- (val);
    \draw[fl] (val)  -- (e1);
    \draw[fl] (mask) |- (bin);
    \draw[fl] (bin)  -- (morf);
    \draw[fl] (morf) -- (dist);
    % El paso de la primera a la segunda hilera rodea por la izquierda.
    \draw[fl] (dist.south) -- ++(0,-4mm) -| (max.north);
    \draw[fl] (max)  -- (ws);
    \draw[fl] (ws)   -- (filt);
    \draw[fl] (filt) -- (e2);

  \end{tikzpicture}
  \end{spacing}
  \caption[Esquema del procedimiento de cálculo]{Esquema del procedimiento. La rama
    superior calcula la cobertura de roca visible (E1) por conteo de píxeles; la inferior
    estima el número de rocas individuales (E2) mediante morfología matemática y división
    de aguas. Las componentes conectadas se cuentan antes de la división de aguas, y una
    componente que no recibe ninguna semilla se cuenta como un único bloque.}
  \label{fig:esquema-metodo}
\end{figure}

\subsection{Cobertura de roca visible (E1)}

Sea $M$ la máscara y $|\cdot|$ el cardinal de un conjunto de píxeles. Se definen el
conjunto de píxeles válidos y el de píxeles de roca como
\[
  V = \{p : M(p) \neq 255\},
  \qquad
  R = \{p : M(p) \in \{1, 3\}\},
\]
es decir, la roca agrupa el lecho rocoso y la roca grande. La cobertura se define como
\begin{equation}
  C = 100 \cdot \frac{|R|}{|V|}.
  \label{eq:cobertura}
\end{equation}

La elección del denominador merece justificación, porque no es neutral. Usar $|V|$
responde a la pregunta \textit{qué fracción del terreno clasificado es roca}, que es la
definición adoptada en la propuesta. La alternativa —dividir por el total de píxeles de la
imagen— responde a una pregunta distinta y entrega una cota inferior conservadora, ya que
trata los píxeles sin etiqueta como si no fueran roca. Ambas medidas se reportan en el
conjunto de resultados, junto con la fracción de escena etiquetada $|V|/N$, donde $N$ es el
número total de píxeles de la imagen, que permite al lector juzgar cuánto pesa la diferencia en
cada escena.

El inconveniente de esta elección es que una escena escasamente etiquetada podría producir
coberturas altas espurias. El Capítulo~\ref{cap:resultados} contrasta explícitamente esa
posibilidad (H1).

\subsection{Conteo aproximado de rocas individuales (E2)}

El conteo opera sobre la máscara binaria $B = \{p : M(p) = 3\}$ y consta de seis etapas.

\paragraph{Limpieza morfológica.} Se aplica una apertura seguida de un cierre, ambos con
elemento estructurante de $3\times3$ píxeles. La apertura elimina píxeles aislados
producto de imprecisiones en el trazado de los polígonos de anotación; el cierre rellena
huecos pequeños dentro de una misma región.

\paragraph{Componentes conectadas.} Se etiquetan las regiones conectadas con criterio de
ocho vecinos. Su número se registra como indicador independiente, porque comparado con el
conteo final permite estimar el efecto neto de la subdivisión y de los filtros.

\paragraph{Transformada de distancia.} Sobre la binaria limpia se calcula la transformada
de distancia euclidiana: cada píxel de roca recibe su distancia al fondo más cercano. La
superficie resultante se suaviza con un filtro gaussiano de desviación estándar
$\sigma = 3{,}0$ píxeles, con el fin de evitar que las irregularidades del contorno
generen máximos espurios.

\paragraph{Detección de semillas por prominencia.} Las semillas de la división de aguas se
obtienen mediante la \textbf{transformada de máximos h}, que conserva únicamente los
máximos cuya altura sobre el entorno supera un umbral $h = 1{,}0$ píxeles de distancia. Se
calcula mediante reconstrucción morfológica \citep{vincent1993reconstruction}. Cuando este
criterio está activo, el parámetro alternativo de separación mínima entre máximos queda sin
efecto en la implementación; por eso no figura entre los parámetros operativos.

Este punto concentra la principal decisión de calibración del trabajo y conviene
explicarlo. La alternativa habitual es detectar máximos locales imponiendo una separación
mínima entre ellos. Con ese criterio, una región extensa de roca grande genera sobre la
transformada de distancia una cresta larga con muchas ondulaciones de poca altura, y cada
una puede dar lugar a una semilla: con la configuración inicial —separación mínima de cinco
píxeles y $\sigma = 1$— una de las escenas de calibración llegó a \cCalSemillasA{} semillas.
El criterio de prominencia descarta las ondulaciones cuya altura sobre el entorno es inferior
a $h$. La Sección~\ref{sec:calibracion} muestra que su efecto no fue selectivo: rebajó el
conteo de forma general, también en escenas sin indicios de sobresegmentación.

\paragraph{División de aguas.} Se inunda el negativo de la transformada de distancia
suavizada partiendo de las semillas, con vecindad de cuatro píxeles y restringiendo el
crecimiento a la máscara binaria \citep{vincent1991watersheds}. El resultado es una partición de la región de roca grande
en subregiones disjuntas. Una componente conectada que no recibe ninguna semilla —por ser
pequeña o casi plana tras el suavizado— se trata como un único bloque, de modo que no se
pierde antes de los filtros, que pueden descartar después tanto fragmentos como componentes
enteras.

\paragraph{Filtrado por tamaño y forma.} Cada subregión debe superar dos criterios para
contarse como roca: un \textbf{área mínima} del \SI{0.05}{\percent} del área de la imagen
—\num{524.3} píxeles para una imagen de \num{1024}$\times$\num{1024}, de modo que la menor
región aceptada tiene \num{525}—, y una
\textbf{relación de aspecto} de su caja envolvente no superior a $5{,}0$. El primero
descarta motas de anotación; el segundo descarta bandas muy alargadas, que en la práctica
corresponden a trazados manuales del horizonte y no a bloques.

\subsection{Parámetros}

Todos los umbrales se pasan como argumentos con valores por defecto documentados; ninguno
está fijado dentro de la lógica de cálculo. La Tabla~\ref{tab:parametros} los recoge.

\begin{table}[htbp]
  \centering
  \caption{Parámetros del procedimiento y su efecto.}
  \label{tab:parametros}
  \begin{tabular}{llp{6.3cm}}
    \toprule
    \textbf{Parámetro} & \textbf{Valor} & \textbf{Efecto} \\
    \midrule
    Apertura            & $3\times3$ px & Elimina motas de anotación. \\
    Cierre              & $3\times3$ px & Rellena huecos internos. \\
    Área mínima         & \SI{0.05}{\percent} & Descarta regiones menores que esa fracción del
      área de la imagen. \\
    Relación de aspecto & $5{,}0$  & Descarta bandas alargadas. \\
    Suavizado $\sigma$  & $3{,}0$  & Evita máximos espurios del contorno. \\
    Prominencia $h$     & $1{,}0$  & Suprime la subdivisión de regiones continuas. \\
    Conectividad        & $8$      & Admite contacto diagonal entre píxeles al etiquetar las
      componentes. \\
    \bottomrule
  \end{tabular}
\end{table}

\subsection{Calibración de los parámetros}
\label{sec:calibracion}

La calibración siguió el esquema: muestra de calibración, parámetros congelados y muestra
independiente de evaluación.

\paragraph{Muestra de calibración.} Seis escenas de MSL NavCam elegidas para cubrir tipos de
terreno distintos: rocas pequeñas dispersas, un cúmulo de bloques próximos, bloques grandes,
una losa continua, un afloramiento de losas fracturadas y una banda alargada trazada sobre el
horizonte. El Anexo~\ref{app:calibracion} las identifica.

\paragraph{Alternativas consideradas.} Se partió de una separación mínima entre máximos de
\SI{5}{\px} con suavizado $\sigma = 1$, que producía sobresegmentación severa —una sola escena
generaba ciento cuarenta y dos semillas, la mayoría a lo largo de una banda alargada—. Se probaron separaciones de \num{15} y
\num{20} píxeles con $\sigma = 3$; la primera preservaba las separaciones reales de los cúmulos
y recortaba los máximos espurios, y la segunda empezaba a fusionar bloques distintos. Persistía
una subdivisión residual de las losas continuas, que se resolvió sustituyendo la separación
mínima por el criterio de prominencia con $h = 1$. La reconstrucción de ese cambio sobre la
población de E2 muestra que \textbf{no fue selectivo}. En las escenas con indicios de
sobresegmentación según la configuración anterior redujo el conteo total en un \cCalRedBaja{}
(solidez media inferior a $0{,}7$, \cCalNBaja{} escenas) y en un \cCalRedMuchas{} (más de quince
rocas, \cCalNMuchas{} escenas), pero también lo redujo en un \cCalRedNorm{} en las \cCalNormN{}
escenas restantes, de las que solo \cCalNormSinCambio{} conservaron su conteo. Una verificación
preliminar, hecha sobre unas pocas escenas, había sugerido que el criterio no alteraba las
escenas normales; la reconstrucción completa no lo sostiene. El criterio de prominencia rebajó el
conteo de forma general, lo que es coherente con el subconteo que revela la validación humana
(Sección~\ref{sec:validacion-manual}).

\paragraph{Reconstrucción.} Un guion del repositorio reconstruye la calibración a partir de un
manifiesto versionado de las seis escenas: evalúa sobre ellas las cuatro configuraciones en el
orden en que se probaron, registra semillas, componentes y rocas, repite la verificación anterior
y comprueba que ninguna escena de calibración pertenece a la muestra de validación. Las cifras de
este apartado proceden de ese guion.

\paragraph{Criterio de elección.} Un parámetro se consideró ajustado cuando el corte reproducía
la separación que un observador trazaría sobre la imagen original. El juicio visual lo hizo el
autor, que es también el observador de la validación. El sesgo esperable de esa coincidencia
favorece el acuerdo entre procedimiento y observador, de modo que, si acaso, atenúa el
resultado negativo de la validación en lugar de producirlo. El área mínima y la relación
de aspecto se fijaron por su función —descartar motas de anotación y bandas del horizonte— y no
se optimizaron.

\paragraph{Congelación e independencia.} Los parámetros quedaron fijados el 31 de agosto de
2026, como consta en el historial del repositorio, antes de recoger ninguna respuesta de la
validación humana. Ninguna de las seis escenas de calibración forma parte de la muestra de
validación definitiva. Cuando la reconstrucción posterior mostró que el criterio de
prominencia no era selectivo (más abajo), se mantuvo $h = 1$: reajustarlo después de la
validación habría supuesto calibrar con la muestra de evaluación.

\paragraph{Sensibilidad.} Para comprobar que las conclusiones no dependen de ese ajuste, se
recalculó el acuerdo con el observador variando tres de los parámetros más influyentes en el
conteo —$h \in \{0{,}5;\,1;\,2;\,3\}$, $\sigma \in \{1;\,3;\,5\}$ y un área mínima de
\SI{0.025}{\percent}, \SI{0.05}{\percent} o \SI{0.1}{\percent}—, \cSenN{} combinaciones en
total. La Sección~\ref{sec:validacion-manual} presenta el resultado.

\subsection{Herramientas y reproducibilidad}

El procedimiento se implementó en Python sobre \texttt{NumPy}
\citep{harris2020numpy}, \texttt{SciPy} \citep{virtanen2020scipy} y
\texttt{scikit-image} \citep{vanderwalt2014skimage}, con \texttt{pandas} para la
construcción del conjunto de resultados. El procedimiento clásico corre íntegramente en
procesador convencional y no requiere unidad gráfica; solo los modelos de aprendizaje
automático de la comparación se ejecutaron sobre la unidad gráfica integrada del equipo. La Tabla~\ref{tab:versiones} registra las versiones
empleadas.

\begin{table}[htbp]
  \centering
  \caption[Versiones de las bibliotecas utilizadas]{Versiones de las bibliotecas utilizadas en la
    ejecución final. El repositorio incluye el registro completo, paquete por paquete.}
  \label{tab:versiones}
  \begin{tabular}{ll}
    \toprule
    \textbf{Componente} & \textbf{Versión} \\
    \midrule
    Python        & 3.13.13 \\
    NumPy         & 2.5.1 \\
    SciPy         & 1.18.0 \\
    scikit-image  & 0.26.0 \\
    pandas        & 3.0.5 \\
    Pillow        & 12.3.0 \\
    Matplotlib    & 3.11.1 \\
    scikit-learn  & 1.9.0 \\
    PyTorch / torchvision & 2.13.0 / 0.28.0 \\
    Ultralytics (FastSAM) & 8.4.137 \\
    \bottomrule
  \end{tabular}
\end{table}

El cálculo de los indicadores sobre las \cNEscenas{} escenas tarda unos diecisiete minutos en
un solo núcleo de un Apple M4 Pro, y unos nueve con ocho procesos en paralelo. Su determinismo
se verificó empíricamente: una reejecución íntegra en un solo núcleo reprodujo de forma idéntica
todas las columnas del conjunto de resultados obtenido en paralelo.

El código se organiza en módulos de responsabilidad única, con funciones que reciben
arreglos y devuelven arreglos o números, sin efectos secundarios. Una única función
concentra la lógica de segmentación y es reutilizada tanto por el cálculo de indicadores
como por la generación de figuras, de modo que las figuras ilustran necesariamente el
mismo cálculo que produce los resultados. El Anexo~\ref{app:recursos} describe la
organización de los recursos.

La reproducibilidad se sostiene además en \textbf{control de versiones}. Todo el material
—código, guiones de ejecución, conjunto de resultados, figuras y la fuente de este
documento— se mantiene en un repositorio Git de historial público
(\url{https://github.com/NotJaayz/mars-agent-research}), de modo que cada
resultado puede rastrearse hasta el estado exacto del código que lo produjo. El conjunto de
datos no se versiona, por su volumen: el código lo referencia mediante una variable de
entorno configurable, y se descarga de la fuente original. Los parámetros no se fijan dentro
de la lógica de cálculo sino que se pasan como argumentos con valores por defecto
documentados. El procedimiento clásico es determinista —no interviene ninguna fuente de
aleatoriedad—, de modo que una reejecución con los mismos parámetros y las mismas versiones
entrega los mismos resultados. Los análisis que sí la emplean —remuestreos, permutaciones,
repartos del segmentador— fijan su semilla.

Tres elementos completan la reproducibilidad. El \textbf{entorno} se declara con las versiones
de la ejecución final, incluidas las bibliotecas de aprendizaje profundo, y se acompaña del
registro exacto de todos los paquetes instalados. Un \textbf{único comando} regenera en orden
los resultados, las evaluaciones, las figuras y todas las cifras de resultados de este
documento, que no se escriben a mano: un guion las extrae de los archivos de resultados y las
inserta en el texto. Por último, \cPruebasN{} \textbf{pruebas automáticas} verifican las reglas deterministas en que descansan
los indicadores —que los píxeles sin etiqueta quedan fuera del denominador, que solo la roca
grande entra en el conteo, que una máscara vacía no tiene cobertura, que una escena casi vacía
no entra en el resumen del conteo, que dos bloques unidos por un estrangulamiento se separan y
que los filtros de área y forma descartan las regiones que no los superan—.

\section{Operacionalización de objetivos}
\label{sec:operacionalizacion}

La Tabla~\ref{tab:operacionalizacion} recoge la matriz de operacionalización, que hace
explícita la correspondencia entre cada objetivo específico, la pregunta que responde, la
técnica empleada, la fuente de información, el resultado esperado y el indicador con que se
verifica. Su función es garantizar la trazabilidad del diseño: ninguna conclusión del
Capítulo~\ref{cap:resultados} carece de un objetivo, una técnica y una métrica asociados.

\begin{landscape}
% Espacio simple y columnas ajustadas al ancho apaisado: con interlineado doble la
% matriz ocupaba tres páginas y quedaba ilegible por el exceso de partición de palabras.
\begin{spacing}{1}
\small
\setlength{\tabcolsep}{4pt}
\begin{longtable}{@{}p{2.3cm}p{3.9cm}p{4.0cm}p{3.3cm}p{4.0cm}p{3.9cm}@{}}
  \caption{Matriz de operacionalización de objetivos, técnicas e indicadores.}
  \label{tab:operacionalizacion}\\
  \toprule
  \textbf{Objetivo} & \textbf{Pregunta} & \textbf{Técnica} & \textbf{Fuente} &
  \textbf{Resultado esperado} & \textbf{Indicador / métrica} \\
  \midrule
  \endfirsthead
  \toprule
  \textbf{Objetivo} & \textbf{Pregunta} & \textbf{Técnica} & \textbf{Fuente} &
  \textbf{Resultado esperado} & \textbf{Indicador / métrica} \\
  \midrule
  \endhead
  \bottomrule
  \endfoot

  \textbf{E1.} Estimar la cobertura de roca visible
  & ¿Qué proporción del terreno clasificado de una escena es roca? (H1)
  & Conteo de píxeles por clase sobre los píxeles etiquetados
  & AI4Mars, máscaras MSL NavCam \textit{train} (secundaria)
  & Indicador de cobertura para toda escena con algún píxel etiquetado
  & \texttt{rock\_coverage\_pct} y \texttt{coverage\_total\_pct}; dependencia con
    \texttt{frac\_valid} (Pearson, Spearman, correlación de distancia, información mutua) \\
  \addlinespace

  \textbf{E2.} Estimar el número de rocas individuales
  & ¿Cuántos bloques discretos se distinguen en las zonas etiquetadas como roca grande? (H2)
  & Morfología matemática, componentes conectadas, transformada de distancia y división de
    aguas con máximos de prominencia
  & Ídem, clase \textit{roca grande}
  & Conteo por imagen y distribución tamaño--frecuencia relativa al campo de visión
  & \texttt{n\_rocks}, \texttt{n\_raw\_components}, \texttt{mean\_solidity}; acuerdo con
    conteo humano (kappa de Cohen simple y ponderado) \\
  \addlinespace

  \textbf{E3.} Caracterizar la composición y la variación
  & ¿Cómo cambian los indicadores según la composición del terreno y a lo largo de la
    secuencia de adquisición? ¿Aportan información distinta? (H4)
  & Estadística descriptiva, tipología de escenas por clase dominante, agregación por
    tramos del reloj de nave
  & Conjunto de resultados propio y metadatos del identificador de imagen
  & Caracterización del terreno y perfil a lo largo de la secuencia
  & \texttt{dominant\_class}, \texttt{scene\_type}, proporción de escenas con roca grande
    por tramo de \texttt{sclk}; dependencia y $\eta^2$ entre los dos indicadores \\
  \addlinespace

  \textbf{E4.} Implementar recursos reproducibles
  & ¿Puede otra persona repetir el análisis y obtener los mismos resultados?
  & Modularización con parámetros explícitos, control de versiones, registro de versiones
    de biblioteca
  & Código propio
  & Archivo tabular con un indicador por imagen, más código y guiones de ejecución
  & \cNEscenas{} filas × 28 columnas en \texttt{results.csv}; reejecución íntegra idéntica \\
  \addlinespace

  \textbf{E5.} Establecer los límites de validez
  & ¿Dependen los indicadores de la fuente de anotación? ¿Reproducen el juicio humano? (H2, H3)
  & Contraste con máscaras de experto, validación con conteo humano, diagnóstico de modos de
    error y comparación con aprendizaje automático
  & Máscaras de experto del propio conjunto (\cNGold); máscaras de rover y distancia;
    conteo de un observador (\cVN{} escenas)
  & Separación de los efectos de imagen y de anotación; atribución de los límites del conteo
  & Descomposición con instrumento común; kappa del conteo con IC; IoU y error con signo del
    segmentador frente a ambas referencias \\

\end{longtable}
\end{spacing}
\end{landscape}

Dos lecturas de esta matriz conviene anticipar. La primera es que los indicadores de E1 y
E2 se verifican con métricas de naturaleza distinta: la cobertura mediante un control
interno de consistencia —no existe verdad de campo con la que compararla— y el conteo
mediante acuerdo con un observador humano. La segunda es que E5 no produce un indicador
propio sino que \textbf{mide la confianza atribuible} a los indicadores anteriores, y de ahí
que sus métricas sean comparaciones entre fuentes y no medidas de terreno.

\section{Estrategia de evaluación y validación}
\label{sec:evaluacion}

Un procedimiento descriptivo como este no admite una métrica de acierto única, porque no existe
verdad de campo sobre la cantidad de roca ni sobre el número de rocas de una escena. La
evaluación se organiza en seis frentes, cuatro de ellos vinculados a las hipótesis H1--H4 de
la Sección~\ref{sec:hipotesis}. Salvo indicación, los contrastes estadísticos emplean intervalos
de confianza por remuestreo \textit{bootstrap} \citep{efron1979bootstrap} y, cuando se trata de
detectar dependencia, valores $p$ por permutación: se barajan los valores de una variable para
romper cualquier relación con la otra y se compara el estadístico observado con esa
distribución nula.

Ambos procedimientos tratan cada escena como una observación intercambiable, y las imágenes no
lo son: se adquieren en secuencias, con una mediana de \cTempEscSolEUno{} escenas por sol en la
población de E1, y la composición del terreno alterna por tramos a lo largo de la misión. Con esa
dependencia los intervalos por escena salen demasiado estrechos y los valores $p$ demasiado
pequeños. Por eso los contrastes principales se repiten con dos procedimientos que la respetan.
El \textbf{\textit{bootstrap} por sol} remuestrea soles completos, con todas sus escenas, en
lugar de escenas sueltas. La \textbf{prueba por desplazamiento circular} ordena las escenas por
reloj de nave y desplaza una variable respecto de la otra: cada serie conserva su
autocorrelación y solo se rompe la relación entre ambas; se excluyen los desplazamientos
menores del \SI{5}{\percent} de la serie.

\paragraph{1. Verificación de la codificación.} Comprobación triple de la correspondencia entre
valores y clases, descrita en la Sección~\ref{sec:datos}. Es la validación más elemental y la
más crítica, porque un error en este punto no produce fallos visibles.

\paragraph{2. Control del denominador (H1).} Si la cobertura dependiera de cuánto se etiquetó,
las escenas escasamente etiquetadas tendrían coberturas espurias. Se contrasta la relación entre
cobertura y fracción etiquetada con medidas de dependencia progresivamente más generales, se
describe la cobertura por tramos de fracción etiquetada y se repite el análisis con la cobertura
sobre la imagen completa como prueba de sensibilidad.

\paragraph{3. Dependencia entre indicadores (H4).} Se mide sobre la población de E2. Un
coeficiente de correlación de Pearson
próximo a cero solo indica ausencia de asociación \textit{lineal}; no demuestra independencia.
Por eso la relación entre cobertura y conteo se mide con un conjunto de estadísticos que
capturan formas de dependencia cada vez más generales: Pearson para la asociación lineal,
Spearman y Kendall para la monótona, la \textbf{correlación de distancia}
\citep{szekely2007dcor} —que vale cero si y solo si las variables son independientes— y la
\textbf{información mutua}, estimada por vecinos próximos \citep{kraskov2004mi}. Se añade una
prueba chi-cuadrado sobre la tabla de contingencia por tramos y la \textbf{razón de correlación}
$\eta^2$, que mide qué fracción de la varianza de un indicador explica el otro sin suponer
ninguna forma funcional. La implementación propia de la correlación de distancia se verificó
antes de usarla sobre casos de respuesta conocida, entre ellos una relación cuadrática que
Pearson no detecta.

\paragraph{4. Contraste entre fuentes de anotación (H3).} Las máscaras colaborativas y las de
especialista cubren imágenes \textbf{distintas}. Una diferencia entre sus coberturas puede
deberse, por tanto, a que cada grupo anote de forma distinta o a que las dos muestras contengan
terreno distinto, y con las etiquetas solas ambos efectos no pueden separarse. Para separarlos se
introduce un \textbf{instrumento común}: el segmentador entrenado, que es una función fija de la
imagen, se aplica a la \textbf{región anotable} de cada escena —todo lo que no es el cuerpo del
rover ni está a más de \SI{30}{\metre}, según las máscaras de rover y de distancia que distribuye
el propio conjunto—, idéntica con independencia de quién anotó. La diferencia de cobertura entre
muestras se descompone entonces en una parte atribuible a las imágenes —la que el instrumento
reproduce sin usar la anotación de las escenas evaluadas— y un residuo atribuible a la
anotación. Como el instrumento se entrenó con etiquetas colaborativas, la descomposición supone
que su error es el mismo en ambas muestras; el Capítulo~\ref{cap:resultados} discute ese
supuesto. La comparación directa de las distribuciones se contrasta con la prueba de
Mann--Whitney y con intervalos para la diferencia de medianas. La muestra colaborativa se toma
al azar entre las imágenes que el segmentador no vio al entrenar, excluidas las de validación;
como puede incluir escenas próximas en el tiempo a imágenes de entrenamiento, se repite como
sensibilidad tomándola solo de los bloques temporales de prueba ($n = \cFPNColab$).

\paragraph{5. Contraste con conteo humano (H2).} Sobre una muestra de \cVN{} escenas de la
población de E2 cuya región de roca grande alcanza al menos \num{2000} píxeles —para que la zona
resaltada sea perceptible— se recogió un conteo manual por bandas —cero, una a tres, cuatro a nueve, diez a
veinticuatro, veinticinco a cuarenta y nueve, cincuenta o más— y se comparó con el automático
mediante el coeficiente kappa de Cohen \citep{cohen1960kappa} y su versión ponderada
\citep{cohen1968weighted} con pesos lineales, apropiada porque las bandas son ordinales; ambos
coeficientes se interpretan según la escala habitual de \citet{landis1977}. La dirección del
error se contrasta con una prueba de signos sobre las escenas en que procedimiento y observador
discrepan. En cada imagen se resaltó la región anotada como roca
grande y se preguntó cuántas rocas se distinguen \textbf{dentro de esa región}: preguntar por
la escena completa habría mezclado dos cuestiones distintas, si la anotación es completa y si el
procedimiento la subdivide bien. Las imágenes se presentaron con identificadores neutros, en
orden barajado y sin el resultado automático. La muestra se estratificó en seis tramos por
cuantiles del área de la región anotada, no por la banda del procedimiento, para no
condicionarla al método evaluado, y excluyó las escenas de la ronda piloto. Una ronda
piloto previa, de \num{24} escenas y cuatro bandas, sirvió para detectar y corregir defectos del
instrumento antes de recoger los datos definitivos. La validación la realizó \textbf{un único
observador}, el propio autor, de modo que no permite separar el desacuerdo atribuible al procedimiento de la
variabilidad propia de la tarea; se presenta como una \textbf{evaluación exploratoria} y no como
una medición de exactitud.

\paragraph{6. Comparación con aprendizaje automático (H2, H3).} Como contraste se emplearon
dos enfoques de aprendizaje profundo. El primero es un modelo fundacional de segmentación
\citep{kirillov2023sam} en su variante rápida \citep{zhao2023fastsam}, aplicado sin
entrenamiento específico a \cSamN{} escenas y restringido a la región de roca grande. El
segundo es un
segmentador binario roca/no-roca, cuya configuración completa recoge la
Tabla~\ref{tab:deeplab-config} para que el experimento sea reproducible.

\begin{spacing}{1}
\small
\begin{longtable}{p{4.0cm}p{10.0cm}}
  \caption[Configuración del segmentador entrenado]{Configuración del segmentador entrenado.}
  \label{tab:deeplab-config}\\
  \toprule
  \textbf{Elemento} & \textbf{Valor} \\
  \midrule
  \endfirsthead
  \toprule
  \textbf{Elemento} & \textbf{Valor} \\
  \midrule
  \endhead
    Arquitectura & DeepLabV3 con red troncal ResNet-50 \citep{chen2017deeplab}, con cabeza
      auxiliar \\
    Pesos iniciales & Preentrenados de \texttt{torchvision} (\texttt{DEFAULT}); ambas cabezas se
      sustituyen por capas de dos clases \\
    Clases & Roca (lecho rocoso y roca grande) y no-roca; los píxeles sin etiqueta se ignoran \\
    Población elegible & Máscaras colaborativas con bandera \texttt{ok}, \texttt{no\_bigrock} o
      \texttt{no\_rock} y fracción etiquetada $\geq 0{,}20$: \cSegElegibles{} escenas. El umbral
      de fracción excluye \cSegExcluidas{} escenas de esas banderas, en las que la pérdida se
      calcularía sobre muy pocos píxeles. Las máscaras de especialista no intervienen \\
    Reparto & Por \textbf{bloques temporales}: las escenas elegibles se ordenan por reloj de nave
      y se dividen en 30 bloques consecutivos de igual tamaño (unos \cSegSolesBloque{} soles cada
      uno), asignados al azar a entrenamiento (20), validación (5) y prueba (5). Se descartan de
      validación y prueba las escenas a menos de un sol de otra partición (\cSegMargen{}) \\
    Muestras & Equilibradas entre escenas con roca (cobertura $> 1$\,\%) y sin roca:
      \cSegTrain{} de entrenamiento, \cSegVal{} de validación y \cSegTest{} de prueba. La prueba
      se repite sobre las \cSegNatN{} escenas elegibles de los bloques de prueba, con su
      distribución natural, y sobre las \cSegExcN{} de esos bloques que excluyó el umbral de
      fracción \\
    Función de pérdida & Entropía cruzada con los píxeles sin etiqueta excluidos, más
      $0{,}4$ veces la de la cabeza auxiliar; sin pesos de clase \\
    Optimizador & Adam, tasa de aprendizaje \cSegLR{} constante, sin planificador \\
    Lote y resolución & \cSegLote{} imágenes de \cSegLado{} $\times$ \cSegLado{} píxeles \\
    Aumento de datos & Ninguno; solo redimensionado y normalización \\
    Épocas y punto de control & \cSegEpocas{} épocas fijas; se conserva el modelo de la época con
      \textbf{mayor mIoU de validación} —media de la intersección sobre la unión de las dos
      clases— (la \cSegEpocaElegida). La regla se fijó antes de
      entrenar y la prueba se evaluó una sola vez, con ese modelo \\
    Decisión & Clase de mayor puntuación por píxel \\
    Semilla & 0 para el reparto y para PyTorch \\
    Equipo y tiempo & Unidad gráfica integrada de un Apple M4 Pro, mediante el motor
      \texttt{MPS} de PyTorch; entre \cSegMinEpocaMin{} y \cSegMinEpocaMax{} minutos por época,
      unos \cSegMinTotal{} en total \\
    Distribución de clases & Roca \cSplitRocaTrain, \cSplitRocaVal{} y \cSplitRocaTest{} de los
      píxeles etiquetados en cada muestra; sin etiqueta, \cSplitIgnTrain, \cSplitIgnVal{} y
      \cSplitIgnTest{} del total \\
    Trazabilidad & Manifiesto con el bloque y la partición de cada escena; huella SHA-256 del
      modelo evaluado (\texttt{\cSegHash}, completa en el Anexo~\ref{app:recursos}). En esa
      unidad gráfica el entrenamiento no es determinista bit a bit aunque se fije la semilla,
      de modo que es la huella, y no la semilla, la que identifica el modelo evaluado \\
    \bottomrule
\end{longtable}
\end{spacing}

El reparto por bloques responde a la posible \textbf{fuga de información}: las imágenes del rover
se adquieren en secuencias, y dos tomas separadas por segundos pueden ser casi idénticas, de modo
que un reparto al azar por imagen deja imágenes casi iguales a ambos lados. Con el margen de un
sol, ninguna imagen de prueba está a menos de \cSegDistMin{} soles de una escena de los bloques
de entrenamiento. La
evaluación distingue además la muestra equilibrada, que mide el desempeño sobre una mezcla
artificial de escenas con y sin roca, de la distribución natural de los bloques de prueba, que
es la que cabe esperar sobre las escenas elegibles. El acuerdo con las etiquetas se mide con la intersección
sobre la unión y, para la cobertura, con el error con signo, su intervalo de confianza y los
límites de acuerdo de \citet{bland1986}. El modelo se evalúa también contra las máscaras de
especialista, que no intervinieron en su entrenamiento ni en su validación.

A estos frentes se añadió, tras observar comportamientos anómalos en el conteo, un
\textbf{diagnóstico sistemático de los modos de error}, que atribuye cada tipo de fallo a su
causa —la etiqueta, los umbrales o el algoritmo—.

\section{Consideraciones éticas y de gestión de datos}

El trabajo no involucra datos personales ni sujetos de investigación: la única participación
humana es el conteo por bandas que realizó el propio autor, sobre imágenes públicas. No requiere,
por tanto, aprobación de un comité de ética. Hay, sin embargo, tres consideraciones pertinentes.

\paragraph{Procedencia y licencia.} Las imágenes son productos públicos de la NASA
obtenidos del Planetary Data System \citep{nasapds2024}, y las máscaras se distribuyen
con el conjunto AI4Mars \citep{swan2022ai4marsdata} bajo licencia Creative Commons
Atribución 4.0. Ambas fuentes
se citan explícitamente. El conjunto de datos, de unos dieciséis gigabytes, no se
redistribuye con el material del proyecto: el código lo referencia por una ruta
configurable, de modo que cualquier persona lo descargue de la fuente original.

\paragraph{Reconocimiento del trabajo de anotación.} Las máscaras existen gracias al
esfuerzo de miles de voluntarios. Un trabajo que examina esa anotación tiene la obligación de
precisar el alcance de sus conclusiones: las diferencias que se documentan en el
Capítulo~\ref{cap:resultados} se refieren a propiedades del conjunto de datos y del diseño de
la tarea de anotación —qué clases existen, cómo se agrega el consenso, qué imágenes integran
cada partición—, no a una deficiencia atribuible a quienes la realizaron.

\paragraph{Alcance de las afirmaciones.} Los indicadores derivados son relativos al
conjunto de máscaras del que provienen y no constituyen estimaciones absolutas de la
abundancia de roca en el terreno marciano. Tampoco están validados para apoyar decisiones de
navegación. Presentarlos como tales sería una sobreinterpretación del dato, y el documento lo
advierte explícitamente en cada resultado afectado.

\section{Cronograma de trabajo}
\label{sec:cronograma}

El trabajo se planificó sobre quince semanas académicas. La
Tabla~\ref{tab:cronograma} resume las actividades previstas por semana.

\begin{longtable}{@{}cp{12.4cm}@{}}
  \caption{Cronograma de trabajo previsto para quince semanas.}
  \label{tab:cronograma}\\
  \toprule
  \textbf{Sem.} & \textbf{Actividades previstas} \\
  \midrule
  \endfirsthead
  \toprule
  \textbf{Sem.} & \textbf{Actividades previstas} \\
  \midrule
  \endhead
  \bottomrule
  \endfoot

  1 & Revisión inicial de la literatura sobre conteo de objetos en imágenes,
      segmentación semántica y métodos de división de aguas. Familiarización con la
      documentación de AI4Mars y con el contexto de operación del rover. \\
  2 & Instalación y prueba del entorno de trabajo: Python, bibliotecas de procesamiento de
      imágenes y control de versiones. Descarga de un subconjunto pequeño para pruebas
      rápidas y comprensión del formato. \\
  3 & Selección del conjunto de imágenes y máscaras definitivo. Organización de los
      directorios de datos y primeros guiones para leer y visualizar las máscaras de forma
      consistente. \\
  4 & Construcción de las dos máscaras binarias de trabajo: una para la cobertura —lecho
      rocoso y roca grande— y otra solo para roca grande. Cálculo preliminar de la
      cobertura por imagen, para verificar que los valores sean razonables. \\
  5 & Implementación del módulo de componentes conectadas sobre la máscara de roca grande.
      Cálculo de área y forma de cada componente y diseño de los filtros iniciales. \\
  6 & Primeros experimentos de conteo usando solo componentes conectadas, sin división de
      aguas. Análisis cualitativo de los casos en que varias rocas quedan unidas en una
      sola región y el conteo resulta subestimado. \\
  7 & Introducción de la transformada de distancia y de un esquema inicial de división de
      aguas para separar rocas en contacto. Pruebas sobre unas pocas imágenes y
      documentación de aciertos y errores. \\
  8 & Ajuste de los parámetros de la división de aguas —suavizado previo, supresión de
      máximos débiles— y observación de la sobresegmentación que aparece en ciertos
      paisajes. Registro de casos problemáticos. \\
  9 & Consolidación del procedimiento completo, desde la máscara original hasta el número
      de rocas por imagen. Correcciones iterativas de parámetros a partir de los ejemplos
      difíciles. \\
  10 & Ejecución del procedimiento sobre el conjunto seleccionado. Revisión de resultados
      atípicos, comparación visual con las imágenes originales y estadísticos descriptivos
      de los indicadores. \\
  11 & Análisis detallado de los resultados e identificación de patrones según el tipo de
      terreno. Notas sobre las limitaciones del método y sobre los casos en que el conteo
      falla con mayor frecuencia. \\
  12 & Redacción de la metodología, con ejemplos gráficos del procedimiento y explicaciones
      dirigidas a lectores no especialistas en visión por computador. \\
  13 & Redacción de resultados y discusión. Integración de las observaciones sobre
      sobresegmentación, sensibilidad a los parámetros y diferencias entre paisajes. \\
  14 & Espacio de maniobra para una actividad opcional de validación con observadores
      humanos —conteos aproximados por bandas— o, de no ser viable, para revisar los
      experimentos y mejorar la presentación de figuras y tablas. \\
  15 & Cierre: revisión completa del documento, verificación de que el código y los datos
      se ejecutan de forma reproducible, armonización del estilo y preparación de la
      versión final. \\
\end{longtable}

Las semanas 1 a 3 se dedican a preparar el terreno: revisar la literatura, dejar listo el
entorno y familiarizarse con las imágenes y máscaras. Entre las semanas 4 y 9 se construye
y ajusta de forma gradual el procedimiento, aceptando que algunos parámetros se afinan por
iteración y no en un único intento. Las semanas 10 a 13 se concentran en aplicarlo sobre el
conjunto elegido, analizar los resultados y plasmarlos en el documento. La semana 14 queda
como margen de maniobra, y la 15 se reserva para el cierre y la verificación de
reproducibilidad.

\subsection*{Ejecución frente a lo planificado}

El plan se cumplió en sus fases principales. Conviene, no obstante, consignar las diferencias
entre lo previsto y lo ejecutado, porque afectan al alcance de los resultados.

\paragraph{La pregunta se precisó.} El anteproyecto preguntaba cómo cuantificar los dos
indicadores. Al contrastarlos con las máscaras de especialista y con el conteo humano se hizo
evidente que la cuestión decisiva era hasta dónde informan sobre el terreno, y la pregunta se
amplió para incluir sus límites de validez. La hipótesis H3 se incorporó en ese momento.

\paragraph{La validación con observadores se realizó con un solo observador.} La semana 14 la
contemplaba como actividad opcional. Se llevó a cabo, y la primera ronda reveló defectos en el
propio instrumento que obligaron a rediseñarlo y repetirla. Se realizó con un único observador,
lo que la limita a una evaluación exploratoria.

\paragraph{Se añadieron tres frentes no previstos.} El contraste con aprendizaje automático se
incorporó al constatar que el procedimiento necesitaba una referencia de contraste, y resultó
decisivo para separar el efecto de la anotación del de las imágenes. El diagnóstico de los
modos de error del conteo surgió al examinar las escenas en que el conteo era nulo. Las reglas heurísticas de
priorización y la aplicación de consulta descritas en el Anexo~\ref{app:extension} surgieron
de la ejecución y no figuraban en el cronograma.
